% In this file you should put the actual content of the blueprint.
% It will be used both by the web and the print version.
% It should *not* include the \begin{document}
%
% If you want to split the blueprint content into several files then
% the current file can be a simple sequence of \input. Otherwise It
% can start with a \section or \chapter for instance.

\section{The Mean Value Theorem}

\begin{definition}
    \label{def:polygonalLine}
    %\lean{PolygonalLine}
    \leanok
    Let $(E, \, ||\cdot||)$ be a normed vector space over $\mathbb{R}$, with subset $A \subseteq E$.
    A \emph{polygonal line} in $A$ is a sequence of $n + 2$ points $u_{0}, \dots, u_{n+1}$ in $E$,
    such that each line segment connecting neighbouring vertices lies completely in $A$, i.e.
    \[
        t \cdot u_{i} + (1 - t) \cdot u_{i + 1} \in A \quad \forall \, i = 0, \ldots, n
        \text{ and } t \in [0,\, 1].
    \]
\end{definition}

\begin{lemma}
    \label{lem:composition_of_polygonal_line}
    %\lean{PolygonalLine.trans}
    \leanok
    \uses{def:polygonalLine}
    Let $U$ be a subset of a normed vector space.
    Given a polygonal line in U, joining $x \in U$ and $y \in U$, together with another polygonal line joining $y$ and $z \in U$,
    then by connecting the two polygonal lines head-to-tail, the result is a polygonal line in U, joining $x$ and $z$.
\end{lemma}
\begin{proof}
    \leanok
    The construction of the composition of polygonal lines follows common sense:
    For the polygonal lines $a$, defined by $(x,\, a_1,\, \dots,\, a_n,\,y)$,
    and $b$, defined by $(y,\, b_1,\, \dots,\, b_m,\, z)$,
    we define the composition $a \circ b$ by the sequence $(x,\,a_1,\,\dots,\, a_n,\,y,\,b_1,\,\dots,\,b_m,\,z)$.

    It remains to show that each segment is still contained in $U$.
    This is trivial as the first $n + 2$ vertices belong to $a$ and are all contained in $U$.
    Similarly, the last $m + 2$ vertices belong to $b$, and are all contained in $U$ as well.
\end{proof}

\begin{theorem}
    \label{thm:existence_of_polygonal_line}
    %\lean{polygonal_connected_of_connected}
    \leanok
    \uses{lem:composition_of_polygonal_line}
    if $U$ is an open connected set of a normed vector space, then any two points of $U$ can be joined by a polygonal line in $U$.
\end{theorem}
\begin{proof}
    \leanok
    The case where $U = \emptyset$ is trivial. So we assume $U$ to be non-empty.
    Fix a point $x \in U$, and consider the set
    \[V \coloneqq \{u \in U \mid \exists \text{ a polygonal line connecting } x \text{ and } u.\}\]
    We want to show that $U = V$.

    To show that, we prove that $V$ is both closed and open in $U$ and $V$ is non-empty.
    The latter is true, as $x$ must be in $V$.

    For closedness, we show that $V = \bar{V}$, the closure of $V$.
    Fix a point $a \in \bar{V}$, then $\exists\, \epsilon > 0$, such that $B_{\epsilon}(a)\,\cap\,V \neq \emptyset$.
    Take a point $b \in B_{\epsilon}(a)\,\cap\,V$.
    Since $b \in V$, there is a polygonal line joining $x$ and $b$.
    Since $b \in B_{\epsilon}(a)$, there is a line segment joining $b$ and $a$, which is a polygonal line as well.
    By composing the two polygonal lines, we create a polygonal line joining $x$ and $a$.
    Hence $a \in V$.

    For openness, fix a point $a \in V$.
    We take an open ball centred at $a$ and contained in $U$, as $U$ is open.
    It can be shown that this open ball is contained in $V$ as well.
    It suffices to show that every point $b$ in the open ball can be connected to $x$ using a polygonal line,
    which can be constructed by composing the polygonal line joining $x$ and $a$ with the line segment joining $a$ and $b$.

    Since $V$ is both open and closed in $U$, it is either equal to $U$ or is the empty set.
    Since it is non-empty, we must have $U = V$.
\end{proof}


\begin{definition}
    \label{def:distance in U}
    %\lean{pathDistance}
    \leanok
    \uses{def:polygonalLine, def:polygonalLineLength}
    Let $U$ be an open connected set of a Banach space $E$. For $a, b \in U$, we define $d_U(a,b)$ to be the infimum of the lengths of the polygonal lines contained in $U$ joining $a$ and $b$.
\end{definition}

\begin{theorem}
    \label{thm:d_U is a distance}
    \leanok
    \uses{def:distance in U}
    $d_U(a,b)$ is a distance in the topological space $U$.
\end{theorem}
\begin{proof}
    \uses{thm:existence_of_polygonal_line, lem:composition_of_polygonal_line, thm:polygonal_line_length_at_least_interval_length, def:polygonalLineLength}
    Let $\varphi$ be an arbitrary polygonal line joining $a$ to $b$ in $U$,
    defined by the sequence $a_0,\, \dots,\, a_{n + 1}$, with $a_0 = a$ and $a_{n + 1} = b$

    \textit{Positivity}:
    The length of $\varphi$ is $\text{Length}(\varphi) := \sum_{i=0}^{n}||a_{i + 1} - a_i||$,
    which by the definition of a norm function is strictly non-negative.
    Since $d_U(a,b)$ is the infimum of a strictly non-negative set, it must itself be non-negative.
    If $d_U(a,b) = 0$, by the previous theorem we have $||b - a|| \leq d_U(a,b) = 0$,
    which implies $a = b$.

    \textit{Symmetry}:
    For any polygonal line $\varphi$ from $a$ to $b$ in $U$,
    the reverse line $\psi$ defined by $a_{n+1},\, \dots,\, a_0$ is a polygonal line from $b$ to $a$.
    $\text{Length}(\psi) = \sum_{i=0}^{n}||a_i - a_{i + 1}|| = \sum_{i=0}^{n}||a_{i+1} - a_i|| = \text{Length}(\varphi)$,
    so the set of lengths of polygonal lines from $a$ to $b$ and from $b$ to $a$ are equivalent and therefore $d_U(a,b) = d_U(b,a)$.

    \textit{Triangle Inequality}:
    Let $\varphi$ be a polygonal line from $a$ to $b$ and $\psi$ be a polygonal line from $b$ to $c$.
    The composition of $\varphi$ and $\psi$ is a polygonal line $\theta$ from $a$ to $c$, where
    \[
        \begin{cases}
            \varphi \text{ is defined by } (a, a_1, \dots, a_n, b) \\
            \psi \text{ is defined by } (b, b_1, \dots, b_m, c)\\
            \theta \text{ is defined by } (a, a_1, \dots, a_n, b, b_1, \dots, b_m, c)
        \end{cases}
    \]
    By the definition of infimum, $d_U(a,c) = \inf \text{Length}(\theta) \leq \text{Length}(\theta) = \text{Length}(\varphi) + \text{Length}(\psi)$.
    Since $\varphi$ and $\psi$ are arbitrary, we take the infimum on the right hand side,
    and hence $d_U(a,c) \leq d_U(a,b) + d_U(b,c)$.
\end{proof}

\begin{theorem}
    \label{thm:mean value theorem in connected open set}
    %\lean{norm_image_sub_le_of_norm_hasFDerivWithin_le'}
    \leanok
    \uses{def:distance in U}
    Let $U$ be a connected open set of a Banach space $E$.
    Let $f: U \rightarrow F$ be a differentiable mapping with values in a Banach space $F$.
    Assume that
    \[
        ||f'(x)|| \leq k \quad \forall x \in U
    \]
    then for any $x_1,x_2 \in U$, $||f(x_1)-f(x_2)|| \leq k \cdot d_U(x_1,x_2)$.
\end{theorem}
\begin{proof}
    \leanok
    \uses{thm:existence_of_polygonal_line, lemma:norm_sum_at_least_sum_norm, def:intervalLength, def:polygonalLineLength, thm:polygonal_line_length_at_least_interval_length}
    Let $x_1,x_2 \in U$ and $l$ a polygonal line contained in $U$ joining $x_1$ and $x_2$.
    Suppose $l$ is defined by $n + 1$ vertices $a_0,\, \dots,\, a_{n+1}$,
    where $a_0=x_1, a_{n+1} = x_2$ and $a_i \in U$.
    Then from the mean value inequality on each interval $(a_i, a_{i+1})$, we have $||f(a_i)-f(a_{i+1})|| \leq k \cdot ||a_i-a_{i+1}||$.
    From the triangle inequality we have
    \begin{align*}
        ||f(x_1)-f(x_2)|| &\leq \sum_{i=0}^{n} ||f(a_i)-f(a_{i+1})|| \\
        &\leq \sum_{i=0}^{n} k \cdot ||a_i-a_{i+1}||\\
        &= k \cdot \text{Length}(l)
    \end{align*}
    as $\text{Length}(l)$ is the sum of the length of each interval.
    Therefore, taking the infimum we get
    \[
        ||f(x_1)-f(x_2)|| \leq k \cdot \inf{\text{Length}(l)} = k\cdot d_U(x_1,x_2)
    \]
\end{proof}

\begin{definition}
    \label{def:intervalLength}
    \leanok
    Let $(E, ||\cdot||)$ be a normed vector space over $\mathbb{R}$.
    The \emph{length of an interval} with origin $a \in E$ and end $b \in E$ is $d(a,b) := ||b - a||$.
\end{definition}

\begin{definition}
    \label{def:polygonalLineLength}
    %\lean{PolygonalLine.length}
    \leanok
    \uses{def:polygonalLine, def:intervalLength}
    Given a polygonal line $\varphi$ defined by $(a_0,\, \dots,\, a_{n+1})$, with some fixed $n \geq 0$
    the \emph{length of the polygonal line} $\varphi$ is the sum of the lengths of the intervals from which it is formed:
    \[
        \text{Length}(\varphi) := \sum_{i=0}^{n} d(a_i, a_{i+1}) = \sum_{i=0}^{n} || a_{i+1} - a_i ||
    \]
\end{definition}

\begin{lemma}
    \label{lemma:norm_sum_at_least_sum_norm}
    %\lean{sublength_sum_geq_distance}
    \leanok
    Given a sequence $(a_i)_{i=0}^{n}$ with $a_i \in E$ for some $n \geq 1$, $\sum_{i=0}^{n-1}||a_{i+1} - a_i|| \geq ||a_n - a_0||$
\end{lemma}
\begin{proof}
    \leanok
    \begin{itemize}
        \item \textbf{Base Case ($n=1$):}
        For $n=1$, we need $||a_1 - a_0|| \geq ||a_1 - a_0||$ which follows immediately by reflexivity.

        \item \textbf{Induction Hypothesis:}
        Assume that the result holds for any such sequence $(b_i)_{i=0}^k$. That is,
        \[
            \sum_{i=0}^{k-1}||b_{i+1} - b_i|| \geq ||b_k - b_0||
        \]

        \item \textbf{Induction Step:}
        We want to show that $\sum_{i=0}^{k}||a_{i+1} - a_i|| \geq ||a_{k+1} - a_0||$ follows from this assumption.
        For some arbitrary sequence  $(a_i)_{i=0}^{k+1}$,
        \begin{align*}
            ||a_{k+1} - a_0|| &=  ||(a_{k+1} - a_k) + (a_k - a_0)|| \\
            &\leq ||a_{k+1} - a_k|| + ||a_k - a_0|| \\
            &\leq ||a_{k+1} - a_k|| + \sum_{i=0}^{k-1}||a_{i+1} - a_i|| \\
            &= \sum_{i=0}^{k}||a_{i+1} - a_i||
        \end{align*}
    \end{itemize}

    Therefore, by the principle of mathematical induction, $\sum_{i=0}^{n-1}||a_{i+1} - a_i|| \geq ||a_n - a_0||$ for any such sequence with any $n \geq 1$.
\end{proof}

\begin{theorem}
    \label{thm:polygonal_line_length_at_least_interval_length}
    \leanok
    % \lean{PolygonalLine.dist_le_length}
    \uses{def:polygonalLineLength, def:intervalLength}
    The length of a polygonal line $\varphi$ defined by $a_0,\, \dots,\, a_{n + 1}$ for some fixed $n \geq 0$
    between origin $a = a_0$ and end $b = a_{n+1}$ is at least equal to the length of the interval between them:
    \[
        \sum_{i=0}^{n}||a_{i+1} - a_i|| \geq ||b - a||
    \]
\end{theorem}
\begin{proof}
    \leanok
    \uses{lemma:norm_sum_at_least_sum_norm}
    The result follows immediately from the lemma.
\end{proof}
